% PWM waveforms at 10% and 90% duty cycle.
%
% Each trace shows three PWM periods of equal length T. The signal is high for
% the first (duty x T) of each period and low for the rest. Pure vector TikZ
% (no opacity/shadings) so it survives the latex -> dvipdf build.
\begin{tikzpicture}[>=Latex, line join=round, font=\footnotesize,
    sig/.style={line width=0.9pt, black},
    ax/.style={->, black!60, line width=0.4pt}]
  \def\T{3}      % period length (cm)
  \def\h{0.9}    % high level (cm)
  \def\gap{2.0}  % vertical spacing between traces (cm)

  % #1 = duty (0-1), #2 = y offset, #3 = label
  \foreach \d/\y/\lbl in {0.1/0/{10\% duty cycle}, 0.9/-\gap/{90\% duty cycle}} {
    \begin{scope}[yshift=\y cm]
      % axes
      \draw[ax] (0,0) -- (3*\T+0.4,0) node[right] {$t$};
      \draw[ax] (0,0) -- (0,\h+0.4);
      \node[left] at (0,\h) {High};
      \node[left] at (0,0) {Low};
      % waveform
      \foreach \k in {0,1,2} {
        \draw[sig] (\k*\T,0) -- (\k*\T,\h) -- (\k*\T+\d*\T,\h)
                   -- (\k*\T+\d*\T,0) -- (\k*\T+\T,0);
      }
      % label
      \node[anchor=west, font=\footnotesize\bfseries] at (3*\T+0.9,0.45*\h) {\lbl};
    \end{scope}
  }

  % period marker on the top trace
  \draw[<->, black!70] (0,\h+0.25) -- node[above] {$T_{\mathrm{PWM}}$} (\T,\h+0.25);
\end{tikzpicture}
